%============================================================
%  ch08_tucker_manifold.tex
%  Chapter 8: The Tucker Manifold
%
%  Tensor Formats in Banach Spaces
%  Antonio Falcó, CEU Universities
%
%  Based on: FHN2019 (Secs. 3-4)
%  Estimated length: ~45 pp.
%============================================================

\chapter{The Tucker Manifold}
\label{chap:tucker_manifold}

This chapter establishes the central geometric result of the monograph:
the set $\Tucker{\fr}{\VD}$ of tensors with fixed Tucker rank carries
the structure of a $\mathcal{C}^\infty$-Banach manifold, and is an
immersed submanifold of the ambient Banach tensor space $\VDnorm$.
Section~\ref{sec:tucker_manifold:decomp} introduces the rank map and
motivates the manifold structure.
Section~\ref{sec:tucker_manifold:charts} constructs the local charts
explicitly. Section~\ref{sec:tucker_manifold:cinfty} proves the
$\mathcal{C}^\infty$ atlas and the fibre bundle structure
(Theorem~\ref{thm:Tucker_Banach_Manifold}).
Section~\ref{sec:tucker_manifold:tangent} identifies the tangent space.
Section~\ref{sec:tucker_manifold:immersion} proves the immersion theorem
(Theorem~\ref{thm:tucker_immersion}).
Throughout, $D = \{1,\ldots,d\}$, $d \geq 2$, and $(V_\alpha,
\|\cdot\|_\alpha)$ are normed spaces.

%------------------------------------------------------------
\section{The rank map and the decomposition of $\VD$}
\label{sec:tucker_manifold:decomp}
%------------------------------------------------------------

Recall from Theorem~\ref{thm:tucker_decomposition_full} that
\begin{equation*}
    \VD = \bigsqcup_{\fr \in \AD{\VD}} \Tucker{\fr}{\VD}.
\end{equation*}
The map
\begin{equation}
\label{eq:rank_map}
    \varrho : \VD \longrightarrow
    \bigcup_{\fr \in \AD{\VD}}
    \prod_{\alpha \in D} \Grass{r_\alpha}{V_\alpha},
    \qquad
    \bv \mapsto \bigl(\Umin{\alpha}{\bv}\bigr)_{\alpha \in D},
\end{equation}
sends each tensor to its tuple of minimal subspaces. For fixed
$\fr \in \AD{\VD}$ with $\fr \neq \mathbf{0}$, the restriction
\begin{equation}
\label{eq:rank_map_restricted}
    \varrho_\fr \coloneqq \varrho\big|_{\Tucker{\fr}{\VD}} :
    \Tucker{\fr}{\VD} \longrightarrow
    \prod_{\alpha \in D} \Grass{r_\alpha}{V_\alpha}
\end{equation}
is surjective: for any tuple $(U_\alpha)_{\alpha \in D}$ with
$U_\alpha \in \Grass{r_\alpha}{V_\alpha}$, any full-rank tensor in
${}_a\bigotimes_\alpha U_\alpha$ belongs to $\Tucker{\fr}{\VD}$ and
maps to $(U_\alpha)$.

The product of Grassmannians $\prod_{\alpha \in D}
\Grass{r_\alpha}{V_\alpha}$ is a Banach manifold (Theorem~\ref{thm:grassmann_manifold}
applied factor by factor), and the rank map $\varrho_\fr$ is its
natural base for the bundle structure of $\Tucker{\fr}{\VD}$.

%------------------------------------------------------------
\section{Local charts}
\label{sec:tucker_manifold:charts}
%------------------------------------------------------------

Fix $\fr \in \AD{\VD}$ with $\fr \neq \mathbf{0}$ and
$\bv \in \Tucker{\fr}{\VD}$. Let $V_{\alpha,\|\cdot\|_\alpha}$
denote the completion of $V_\alpha$ and choose for each $\alpha \in D$
a closed complement $\Wmin{\alpha}{\bv}$ of $\Umin{\alpha}{\bv}$ in
$V_{\alpha,\|\cdot\|_\alpha}$:
\begin{equation*}
    V_{\alpha,\|\cdot\|_\alpha}
    = \Umin{\alpha}{\bv} \oplus \Wmin{\alpha}{\bv}.
\end{equation*}
Such a complement exists by Proposition~\ref{prop:finitedim_split},
since $\Umin{\alpha}{\bv}$ is finite-dimensional.

\subsection*{Chart domain}

Define the \emph{chart domain at $\bv$} by
\begin{equation}
\label{eq:chart_domain}
    \chart(\bv) \coloneqq
    \bigl\{
        \bw \in \Tucker{\fr}{\VD} :
        \Umin{\alpha}{\bw} \in
        \GrassB{\Wmin{\alpha}{\bv}, V_\alpha},\;
        \forall \alpha \in D
    \bigr\},
\end{equation}
where $\GrassB{W, X} = \{U \in \GrassB{X} : X = U \oplus W\}$ is
the chart domain of the Grassmannian at $\Umin{\alpha}{\bv}$
(see Section~\ref{sec:banach:grassmann}).

\subsection*{Chart map}

Fix bases $\{u^{(\alpha)}_{i_\alpha}\}_{i_\alpha=1}^{r_\alpha}$ of
$\Umin{\alpha}{\bv}$ for each $\alpha \in D$. By
Lemma~\ref{lem:tucker_characterisation}, there is a unique
$\CD(\bv) \in \RR^{\times r_\alpha}_*$ such that $\bv$ has the
Tucker representation~\eqref{eq:tucker_core}.
For each $\bw \in \chart(\bv)$, define the \emph{chart map}
\begin{equation}
\label{eq:chart_map}
    \chartmap{\bv} : \chart(\bv)
    \longrightarrow
    \prod_{\alpha \in D} \cL\bigl(\Umin{\alpha}{\bv},
        \Wmin{\alpha}{\bv}\bigr)
    \times \RR^{\times r_\alpha}_*,
\end{equation}
by sending $\bw = (\bigotimes_\alpha \exp(L_\alpha))(\bu(E^{(D)}))$
to the pair $((L_\alpha)_{\alpha \in D},\, E^{(D)})$, where
$(L_\alpha, E^{(D)})$ are uniquely determined by
Lemma~\ref{lem:element} below.

\begin{lemma}[Local representation]
\label{lem:element}
Under the conditions above, for $\bw \in \Tucker{\fr}{\VD}$
the following are equivalent:
\begin{enumerate}
\item[\textup{(a)}] $\bw \in \chart(\bv)$.
\item[\textup{(b)}] There exists a unique pair
    $\bigl((L_\alpha)_{\alpha \in D},\, E^{(D)}\bigr)
    \in \prod_\alpha \cL\bigl(\Umin{\alpha}{\bv}, \Wmin{\alpha}{\bv}\bigr)
    \times \RR^{\times r_\alpha}_*$
    such that
    \begin{equation}
    \label{eq:local_rep}
        \bw
        = \Bigl(\bigotimes_{\alpha \in D} \exp(L_\alpha)\Bigr)
          \bigl(\bu(E^{(D)})\bigr),
    \end{equation}
    where $\bu(E^{(D)}) \in \Tucker{\fr}\bigl({}_a\bigotimes_\alpha
    \Umin{\alpha}{\bv}\bigr)$ is the tensor with core $E^{(D)}$ in
    the fixed bases $\{u^{(\alpha)}_{i_\alpha}\}$.
\end{enumerate}
\end{lemma}

\begin{proof}
We prove (a)$\Rightarrow$(b) and (b)$\Rightarrow$(a).

\medskip
\noindent\textbf{(a) $\Rightarrow$ (b): Existence and uniqueness.}

Suppose $\bw \in \chart(\bv)$, so $\Umin{\alpha}{\bw} \in
\GrassB{\Wmin{\alpha}{\bv}, V_\alpha}$ for every $\alpha \in D$.

\textit{Step 1: Recovering $L_\alpha$.}
Define $L_\alpha$ as the Grassmannian chart coordinate of
$\Umin{\alpha}{\bw}$ at $\Umin{\alpha}{\bv}$:
\begin{equation*}
    L_\alpha
    \coloneqq \Psi_{\Umin{\alpha}{\bv} \oplus \Wmin{\alpha}{\bv}}
    \bigl(\Umin{\alpha}{\bw}\bigr)
    = \Proj{\Wmin{\alpha}{\bv}}{\Umin{\alpha}{\bv}}
      \big|_{\Umin{\alpha}{\bw}}
      \circ
      \bigl(\Proj{\Umin{\alpha}{\bv}}{\Wmin{\alpha}{\bv}}
      \big|_{\Umin{\alpha}{\bw}}\bigr)^{-1}
    \in \cL\bigl(\Umin{\alpha}{\bv}, \Wmin{\alpha}{\bv}\bigr).
\end{equation*}
By Proposition~\ref{prop:nilpotent_algebra},
$\exp(L_\alpha) = \Id{V_\alpha} + L_\alpha$, and
\begin{equation}
\label{eq:exp_maps_U_to_Uw}
    \exp(L_\alpha)\bigl(\Umin{\alpha}{\bv}\bigr)
    = G(L_\alpha)
    = \Umin{\alpha}{\bw},
\end{equation}
since $G(L_\alpha)$ is the Grassmann chart inverse of $L_\alpha$
(Theorem~\ref{thm:grassmann_manifold}).

\textit{Step 2: Recovering $E^{(D)}$.}
Set $\widetilde{\bw} \coloneqq
\bigl(\bigotimes_{\alpha \in D} \exp(-L_\alpha)\bigr)(\bw)$.
Since $\exp(-L_\alpha) = \Id{} - L_\alpha = (\exp(L_\alpha))^{-1}$
maps $\Umin{\alpha}{\bw}$ back to $\Umin{\alpha}{\bv}$
by~\eqref{eq:exp_maps_U_to_Uw}, we get
$\Umin{\alpha}{\widetilde\bw} \subset \Umin{\alpha}{\bv}$ for each
$\alpha$. Proposition~\ref{prop:min_monotone} gives
$\widetilde\bw \in {}_a\bigotimes_\alpha \Umin{\alpha}{\bv}$,
and since $\exp(-L_\alpha)$ preserves ranks,
$\widetilde\bw \in \Tucker{\fr}({}_a\bigotimes_\alpha \Umin{\alpha}{\bv})$.
Its core array in the fixed bases is the unique $E^{(D)} \in
\RR^{\times r_\alpha}_*$ such that $\bu(E^{(D)}) = \widetilde\bw$.
Applying $\bigotimes_\alpha \exp(L_\alpha)$ recovers~\eqref{eq:local_rep}.

\textit{Uniqueness.}
If $(L_\alpha, E^{(D)})$ and $(L'_\alpha, E'^{(D)})$ both represent $\bw$,
applying $\bigotimes_\alpha \exp(-L'_\alpha)$ and Step~2 gives
$E^{(D)} = E'^{(D)}$, and then $L_\alpha = L'_\alpha$ for each $\alpha$.

\medskip
\noindent\textbf{(b) $\Rightarrow$ (a).}

For $(L_\alpha, E^{(D)})$ with $E^{(D)} \in \RR^{\times r_\alpha}_*$,
let $\bw$ be defined by~\eqref{eq:local_rep}. By~\eqref{eq:exp_maps_U_to_Uw},
$\Umin{\alpha}{\bw} = G(L_\alpha) \in \GrassB{\Wmin{\alpha}{\bv}, V_\alpha}$,
so $\bw \in \chart(\bv)$.
\end{proof}

Here $\exp(L_\alpha) = \Id{V_\alpha} + L_\alpha$ by the nilpotency
of $L_\alpha \in \cL_{(\Umin{\alpha}{\bv},\Wmin{\alpha}{\bv})}(V_\alpha)$
(Proposition~\ref{prop:nilpotent_algebra}), so the representation is
polynomial and $\mathcal{C}^\infty$.

\begin{remark}
\label{rem:chart_geometric}
Formula~\eqref{eq:local_rep} has a clean geometric interpretation:
$\bw$ is obtained from the ``reference'' tensor $\bu(E^{(D)})$
(which lives in the fixed subspaces $\Umin{\alpha}{\bv}$) by applying
the ``rotation'' $\bigotimes_\alpha \exp(L_\alpha)$ which moves each
subspace $\Umin{\alpha}{\bv}$ to $\Umin{\alpha}{\bw} = G(L_\alpha)
= \spann\{(\Id{} + L_\alpha)(u^{(\alpha)}_{i_\alpha})\}$.
\end{remark}

%------------------------------------------------------------
\section{The $\mathcal{C}^\infty$-Banach manifold structure}
\label{sec:tucker_manifold:cinfty}
%------------------------------------------------------------

\begin{lemma}[Smooth transitions]
\label{lem:premanifold}
Under the continuity assumption on the tensor product map
\eqref{eq:tensor_product_normed}, for $\bv, \bv' \in
\Tucker{\fr}{\VD}$ with $\chart(\bv) \cap \chart(\bv') \neq \emptyset$,
the transition map
\begin{equation*}
    \chartmap{\bv'} \circ \chartmap{\bv}^{-1} :
    \chartmap{\bv}\bigl(\chart(\bv) \cap \chart(\bv')\bigr)
    \longrightarrow
    \chartmap{\bv'}\bigl(\chart(\bv) \cap \chart(\bv')\bigr)
\end{equation*}
is $\mathcal{C}^\infty$-Fréchet differentiable \textup{(}analytic if
the $V_\alpha$ are complex Banach spaces\textup{)}.
\end{lemma}

\begin{proof}
Let $\bw \in \chart(\bv) \cap \chart(\bv')$ with chart coordinates
$\chartmap{\bv}(\bw) = ((L_\alpha), E^{(D)})$ and
$\chartmap{\bv'}(\bw) = ((L'_\alpha), E'^{(D)})$.
We show $(L_\alpha, E^{(D)}) \mapsto (L'_\alpha, E'^{(D)})$ is
$\mathcal{C}^\infty$ (analytic over $\CC$).

\medskip
\noindent\textbf{Step 1: Smoothness of the Grassmannian component.}

For each $\alpha \in D$, $\Umin{\alpha}{\bw} = G(L_\alpha)$
(the inverse Grassmann chart of $L_\alpha$ at $\Umin{\alpha}{\bv}$).
The map
\begin{equation*}
    L_\alpha \mapsto L'_\alpha
    = \Psi_{\Umin{\alpha}{\bv'} \oplus \Wmin{\alpha}{\bv'}}(G(L_\alpha))
\end{equation*}
is the transition map of $\Grass{r_\alpha}{V_\alpha}$ between the
charts at $\Umin{\alpha}{\bv}$ and $\Umin{\alpha}{\bv'}$. Since
$\Grass{r_\alpha}{V_\alpha}$ is an analytic Banach manifold
(Theorem~\ref{thm:grassmann_manifold}), its transition maps are
analytic, hence $L_\alpha \mapsto L'_\alpha$ is analytic.
More explicitly: $G(L_\alpha) = \{(\Id{}+L_\alpha)(u) : u \in \Umin{\alpha}{\bv}\}$
is affine-linear in $L_\alpha$; the projections in the chart formula
\eqref{eq:grassmann_chart} are bounded linear; and inversion of a
bounded linear isomorphism is analytic (locally given by the Neumann
series). Their composition is analytic.

\medskip
\noindent\textbf{Step 2: Smoothness of the core component.}

Given $(L_\alpha)$ (and hence $\Umin{\alpha}{\bw} = G(L_\alpha)$),
the core is determined by
\begin{equation*}
    \bu(E^{(D)})
    = \Bigl(\bigotimes_{\alpha \in D} \exp(-L_\alpha)\Bigr)(\bw),
\end{equation*}
where $\exp(-L_\alpha) = \Id{} - L_\alpha$
(Proposition~\ref{prop:nilpotent_algebra}).
The map $(L_\alpha, \bw) \mapsto (\bigotimes_\alpha (\Id{}-L_\alpha))(\bw)$
is multilinear and continuous in $(L_\alpha)$ and $\bw$, hence
$\mathcal{C}^\infty$ by Proposition~\ref{prop:multilinear_smooth}
(analytic over $\CC$ by Corollary~\ref{cor:multilinear_analytic}).
Reading off the core array $E^{(D)}$ from the result is a bounded
linear operation. The same applies to $E'^{(D)}$ using $(L'_\alpha)$
from Step~1.

\medskip
\noindent\textbf{Conclusion.}

The transition map $(L_\alpha, E^{(D)}) \mapsto (L'_\alpha, E'^{(D)})$
is a composition of $\mathcal{C}^\infty$ (analytic over $\CC$) maps,
hence $\mathcal{C}^\infty$ (analytic over $\CC$).
\end{proof}

\begin{theorem}[$\mathcal{C}^\infty$ Tucker manifold,
{\cite[Theorem~3.17]{FHN2019}}]
\label{thm:Tucker_Banach_Manifold}
Let $(V_\alpha, \|\cdot\|_\alpha)$ be normed spaces for $\alpha \in D$
and let $\|\cdot\|_D$ be a norm on $\VD$ such that the tensor product
map~\eqref{eq:tensor_product_normed} is continuous. Then:
\begin{enumerate}
\item[\textup{(i)}] The collection $\{(\chart(\bv), \chartmap{\bv})\}_{\bv \in
    \Tucker{\fr}{\VD}}$ is a $\mathcal{C}^\infty$-atlas on
    $\Tucker{\fr}{\VD}$, which is thus a $\mathcal{C}^\infty$-Banach
    manifold modelled on the Banach space
    \begin{equation}
    \label{eq:tucker_model_space}
        \Emodel{D} \coloneqq
        \prod_{\alpha \in D}
        \cL\bigl(\Umin{\alpha}{\bv}, \Wmin{\alpha}{\bv}\bigr)
        \times \RR^{\times_{\alpha \in D} r_\alpha}.
    \end{equation}
\item[\textup{(ii)}] The triple $(\Tucker{\fr}{\VD},\,
    \prod_{\alpha \in D}\Grass{r_\alpha}{V_\alpha},\, \varrho_\fr)$
    is a $\mathcal{C}^\infty$-fibre bundle with typical fibre
    $\RR^{\times r_\alpha}_*$.
\end{enumerate}
If the $V_\alpha$ are complex Banach spaces, the manifold is analytic.
\end{theorem}

\begin{proof}
The atlas conditions AT1--AT3 are verified as follows.
AT1: $\{\chart(\bv)\}$ covers $\Tucker{\fr}{\VD}$ since every
$\bv$ belongs to its own chart domain.
AT2: $\chartmap{\bv}$ is a bijection onto
$\prod_\alpha \cL(\Umin{\alpha}{\bv}, \Wmin{\alpha}{\bv}) \times
\RR^{\times r_\alpha}_*$ by Lemma~\ref{lem:element}.
AT3: Lemma~\ref{lem:premanifold} gives $\mathcal{C}^\infty$ transitions.
The fibre bundle structure follows because $\varrho_\fr$ acts locally
as the projection onto $\prod_\alpha \GrassB{\Wmin{\alpha}{\bv}, V_\alpha}$
(the Grassmannian chart domain), with fibre $\RR^{\times r_\alpha}_*$.
\end{proof}

\begin{remark}[Independence from the norm]
\label{rem:norm_independence}
The manifold structure of $\Tucker{\fr}{\VD}$ is independent of the
choice of $\|\cdot\|_D$: the chart maps $\chartmap{\bv}$ depend only
on the splitting $V_{\alpha,\|\cdot\|_\alpha} = \Umin{\alpha}{\bv}
\oplus \Wmin{\alpha}{\bv}$, and the $\mathcal{C}^\infty$ regularity
follows from the continuity of the tensor product map, which holds for
any crossnorm. Different norms give equivalent smooth structures.
\end{remark}

\begin{corollary}[Hilbert manifold]
\label{cor:hilbert_manifold}
If $V_{\alpha,\|\cdot\|_\alpha}$ is a Hilbert space for each
$\alpha \in D$, then $\Tucker{\fr}{\VD}$ is a $\mathcal{C}^\infty$-Hilbert
manifold modelled on
$\prod_\alpha W^{r_\alpha}_\alpha \times \RR^{\times r_\alpha}$,
where $W^{r_\alpha}_\alpha = (U_\alpha)^\perp \times \cdots$ (each
$\cL(U_\alpha, W_\alpha) \cong W_\alpha^{r_\alpha}$ via the chosen
orthonormal basis of $U_\alpha$).
\end{corollary}

\begin{example}
\label{ex:tucker_mfd_Lp}
For $1 < p < \infty$, the $L^p$ tensor space
$L^p(\prod_\alpha I_\alpha)$ satisfies the hypotheses
of Theorem~\ref{thm:Tucker_Banach_Manifold}. Hence
$\Tucker{\fr}({}_a\bigotimes_\alpha L^p(I_\alpha))$ is a
$\mathcal{C}^\infty$-Banach manifold. For $p=2$ it is a Hilbert
manifold.
\end{example}

\begin{example}
\label{ex:tucker_mfd_Sobolev}
For $V_1 = H^{1,p}(I_1)$, $V_2 = H^{1,p}(I_2)$, and either the
full Sobolev norm $\|\cdot\|_{H^{1,p}}$ or the partial norm
$\|\cdot\|_{(0,1),p}$ on $V_1 \atens V_2$, both norms make the
tensor product map continuous (see~\cite[Examples~4.41,~4.42]{Hackbusch2019}).
Hence $\Tucker{(r,r)}{H^{1,p}(I_1) \atens H^{1,p}(I_2)}$ is a
$\mathcal{C}^\infty$-Banach manifold modelled on
$\cL(U_1, W_1) \times \cL(U_2, W_2) \times \GL{\RR^r}$.
\end{example}

%------------------------------------------------------------
\section{The tangent space}
\label{sec:tucker_manifold:tangent}
%------------------------------------------------------------

\begin{definition}[Tangent subspace $Z^{(D)}(\bv)$]
\label{def:ZD}
For $\bv \in \Tucker{\fr}{\VD}$, define the linear subspace
\begin{equation}
\label{eq:ZD}
    \ZD{\bv} \coloneqq
    {}_a\bigotimes_{\alpha \in D} \Umin{\alpha}{\bv}
    \oplus
    \bigoplus_{\alpha \in D}
    \Wmin{\alpha}{\bv} \atens \Umin{D\setminus\{\alpha\}}{\bv}
    \;\subset\; \VDnorm.
\end{equation}
\end{definition}

$\ZD{\bv}$ is a direct sum of $d+1$ subspaces: the core subspace
${}_a\bigotimes_\alpha \Umin{\alpha}{\bv}$ (dimension $\prod r_\alpha$)
and $d$ ``tangent directions'' $\Wmin{\alpha}{\bv} \atens
\Umin{D\setminus\{\alpha\}}{\bv}$.

\begin{proposition}[Tangent space identification]
\label{prop:tangent_space}
Under the continuity assumption, the tangent map of the inclusion
$\mathfrak{i}: \Tucker{\fr}{\VD} \hookrightarrow \VDnorm$ at $\bv$ is
injective, and
\begin{equation*}
    T_\bv \mathfrak{i}\bigl(\mathbb{T}_\bv(\Tucker{\fr}{\VD})\bigr)
    = \ZD{\bv}.
\end{equation*}
More explicitly, for $(\dot{\mathfrak{L}}, \dot{C}^{(D)}) \in
\mathbb{T}_\bv(\Tucker{\fr}{\VD})$:
\begin{equation}
\label{eq:tangent_formula}
    T_\bv \mathfrak{i}(\dot{\mathfrak{L}}, \dot{C}^{(D)})
    = \sum_{(i_\alpha)} \dot{C}^{(D)}_{(i_\alpha)}
      \bigotimes_\alpha u^{(\alpha)}_{i_\alpha}
    + \sum_{\alpha \in D} \sum_{i_\alpha=1}^{r_\alpha}
      \dot{u}^{(\alpha)}_{i_\alpha} \otimes \bU^{(\alpha)}_{i_\alpha},
\end{equation}
where $\dot{u}^{(\alpha)}_{i_\alpha} = \dot{L}_\alpha(u^{(\alpha)}_{i_\alpha})
\in \Wmin{\alpha}{\bv}$ and $\bU^{(\alpha)}_{i_\alpha} \in
\Umin{D\setminus\{\alpha\}}{\bv}$ is given by~\eqref{eq:tucker_Ualpha}.
\end{proposition}

\begin{proof}
We compute $T_\bv \mathfrak{i}$ by differentiating
$\mathfrak{i} \circ \chartmap{\bv}^{-1}$ at the base point
$(0, C_\bv^{(D)}) = \chartmap{\bv}(\bv)$.

\medskip
\noindent\textbf{Step 1: The inverse chart map is degree-one polynomial.}

By Lemma~\ref{lem:element} and Proposition~\ref{prop:nilpotent_algebra},
$\exp(L_\alpha) = \Id{} + L_\alpha$, so:
\begin{align}
\label{eq:chart_inv_poly}
    (\mathfrak{i} \circ \chartmap{\bv}^{-1})(L_\alpha, E^{(D)})
    &= \Bigl(\bigotimes_\alpha (\Id{} + L_\alpha)\Bigr)
       \Bigl(\sum_{(i_\alpha)} E^{(D)}_{(i_\alpha)}
       \bigotimes_\alpha u^{(\alpha)}_{i_\alpha}\Bigr).
\end{align}
Expanding $\bigotimes_\alpha(\Id{}+L_\alpha)$ by multilinearity and
using the nilpotency $L_\alpha \circ L_\beta = 0$ for all $\alpha, \beta$
(Proposition~\ref{prop:nilpotent_algebra}), all terms of degree $\geq 2$
in $(L_\alpha)$ vanish:
\begin{equation*}
    \bigotimes_\alpha (\Id{} + L_\alpha)
    = \bigotimes_\alpha \Id{}
    + \sum_{\alpha \in D}
      L_\alpha \otimes \bigotimes_{\beta \neq \alpha} \Id{\beta}
    \pmod{\text{degree} \geq 2}.
\end{equation*}
Hence~\eqref{eq:chart_inv_poly} is exactly affine-linear in
$(L_\alpha, E^{(D)})$.

\medskip
\noindent\textbf{Step 2: Computing the Fréchet derivative.}

Differentiating~\eqref{eq:chart_inv_poly} at $(0, C_\bv^{(D)})$ in
direction $(\dot{L}_\alpha, \dot{C}^{(D)})$:
\begin{align*}
    D(\mathfrak{i} \circ \chartmap{\bv}^{-1})\big|_{(0,C_\bv^{(D)})}
    (\dot{L}_\alpha, \dot{C}^{(D)})
    &= \sum_{(i_\alpha)} \dot{C}^{(D)}_{(i_\alpha)}
       \bigotimes_\alpha u^{(\alpha)}_{i_\alpha}\\
    &\quad+ \sum_{\alpha \in D}
       \sum_{(i_\alpha)} C^{(D)}_{\bv,(i_\alpha)}\,
       \dot{L}_\alpha(u^{(\alpha)}_{i_\alpha})
       \otimes \bigotimes_{\beta \neq \alpha} u^{(\beta)}_{i_\beta}.
\end{align*}
Using the definition of $\bU^{(\alpha)}_{i_\alpha}$
(formula~\eqref{eq:tucker_Ualpha}), the second line equals
$\sum_\alpha \sum_{i_\alpha}
\dot{L}_\alpha(u^{(\alpha)}_{i_\alpha}) \otimes \bU^{(\alpha)}_{i_\alpha}$,
giving formula~\eqref{eq:tangent_formula} with
$\dot{u}^{(\alpha)}_{i_\alpha} = \dot{L}_\alpha(u^{(\alpha)}_{i_\alpha})
\in \Wmin{\alpha}{\bv}$.

\medskip
\noindent\textbf{Step 3: Image equals $\ZD{\bv}$, and injectivity.}

Formula~\eqref{eq:tangent_formula} shows $T_\bv\mathfrak{i}$ maps into
$\ZD{\bv}$: the first term lies in
${}_a\bigotimes_\alpha \Umin{\alpha}{\bv}$ and the $\alpha$-th summand
in the second term lies in
$\Wmin{\alpha}{\bv} \atens \Umin{D\setminus\{\alpha\}}{\bv}$.

For surjectivity: given $\bz \in \ZD{\bv}$ decomposed as
$\bz = \bz_0 + \sum_\alpha \bz_\alpha$ (cf.~Definition~\ref{def:ZD}),
write $\bz_0 = \sum_{(i_\alpha)} \dot{C}^{(D)}_{(i_\alpha)}
\bigotimes_\alpha u^{(\alpha)}_{i_\alpha}$ to determine $\dot{C}^{(D)}$,
and $\bz_\alpha = \sum_{i_\alpha} w^{(\alpha)}_{i_\alpha} \otimes
\bU^{(\alpha)}_{i_\alpha}$ with $w^{(\alpha)}_{i_\alpha} \in
\Wmin{\alpha}{\bv}$, defining $\dot{L}_\alpha$ by
$\dot{L}_\alpha(u^{(\alpha)}_{i_\alpha}) = w^{(\alpha)}_{i_\alpha}$.
Then $T_\bv\mathfrak{i}(\dot{L}_\alpha, \dot{C}^{(D)}) = \bz$.

Injectivity: if $T_\bv\mathfrak{i}(\dot{L}_\alpha, \dot{C}^{(D)}) = 0$,
then $\bz_0 = 0$ gives $\dot{C}^{(D)} = 0$, and $\bz_\alpha = 0$
for each $\alpha$ gives $\dot{L}_\alpha = 0$.
\end{proof}

%------------------------------------------------------------
\section{Immersion in the ambient Banach space}
\label{sec:tucker_manifold:immersion}
%------------------------------------------------------------

To show that $\Tucker{\fr}{\VD}$ is an immersed submanifold of
$\VDnorm$, it remains to show that $\ZD{\bv} \in \GrassB{\VDnorm}$,
i.e., that $\ZD{\bv}$ is a split subspace of the ambient Banach space.
This requires condition~\condINJ.

\begin{lemma}
\label{lem:two_faces}
Assume \condINJ\ holds. Let $\beta \in D$ and $W_\beta \in
\GrassB{V_{\beta,\|\cdot\|_\beta}}$ with $V_{\beta,\|\cdot\|_\beta}
= U_\beta \oplus W_\beta$ for some finite-dimensional $U_\beta$.
Then for every finite-dimensional subspace $U_{[\beta]} \subset
\Vbigalpha^{[\beta]}$:
\begin{equation*}
    W_\beta \atens U_{[\beta]} \in \GrassB{\VDnorm}.
\end{equation*}
\end{lemma}

\begin{proof}
We construct an explicit bounded idempotent with image
$W_\beta \atens U_{[\beta]}$.

\medskip
\noindent\textbf{Step 1: Bounded projection onto $U_{[\beta]}$.}

Since $U_{[\beta]}$ is finite-dimensional,
Proposition~\ref{prop:finitedim_split} provides a complement
$W_{[\beta]}$ with $\Vbigalpha^{[\beta]} = U_{[\beta]} \oplus W_{[\beta]}$.
Let $Q_{[\beta]} \coloneqq \Proj{U_{[\beta]}}{W_{[\beta]}}$.
The map $\overline{\Id{\beta} \otimes Q_{[\beta]}} \in \cL(\VDnorm)$
is well-defined and bounded: on elementary tensors,
$\|v_\beta \otimes Q_{[\beta]}(\bU)\|_D = \|v_\beta\|_\beta
\|Q_{[\beta]}(\bU)\|_{[\beta]} \leq \|v_\beta\|_\beta
\|Q_{[\beta]}\|\,\|\bU\|_{[\beta]}$
using \condCROSS, and the extension to $\VDnorm$ by continuity is
bounded with norm $\leq \|Q_{[\beta]}\|$.

\medskip
\noindent\textbf{Step 2: Construction of $\mathcal{P}$.}

Let $P_\beta \coloneqq \Proj{W_\beta}{U_\beta} \in
\cL(V_{\beta,\|\cdot\|_\beta})$ be bounded (since $W_\beta$ is split).
Define
\begin{equation*}
    \mathcal{P}
    \coloneqq \overline{P_\beta \otimes \Id{[\beta]}}
    \circ
    \overline{\Id{\beta} \otimes Q_{[\beta]}}
    \in \cL(\VDnorm).
\end{equation*}
This is a composition of bounded operators, hence bounded.

\textit{Idempotency.} On elementary tensors:
$\mathcal{P}^2(v_\beta \otimes \bU)
= P_\beta^2(v_\beta) \otimes Q_{[\beta]}^2(\bU)
= P_\beta(v_\beta) \otimes Q_{[\beta]}(\bU)
= \mathcal{P}(v_\beta \otimes \bU)$.
By density and continuity, $\mathcal{P}^2 = \mathcal{P}$.

\textit{Image.} On elementary tensors:
$\mathcal{P}(v_\beta \otimes \bU)
= P_\beta(v_\beta) \otimes Q_{[\beta]}(\bU)
\in W_\beta \atens U_{[\beta]}$.
Conversely, for $w_\beta \in W_\beta$ and $u \in U_{[\beta]}$:
$\mathcal{P}(w_\beta \otimes u) = w_\beta \otimes u$,
since $P_\beta w_\beta = w_\beta$ and $Q_{[\beta]} u = u$.
Hence $\mathrm{Im}\,\mathcal{P} = W_\beta \atens U_{[\beta]}$.

\medskip
By Proposition~\ref{prop:split_equiv}(b),
$W_\beta \atens U_{[\beta]} \in \GrassB{\VDnorm}$.
\end{proof}

\begin{lemma}[{\cite[Lemma~4.14]{FHN2019}}]
\label{lem:direct_sum_G}
Let $X$ be a Banach space and $U, V \in \GrassB{X}$ with
$U \cap V = \{0\}$. Then $U \oplus V \in \GrassB{X}$ and
$U \cap V \in \GrassB{X}$.
\end{lemma}

\begin{proof}
Since $U, V \in \GrassB{X}$ there exist $U', V' \in \GrassB{X}$ with
$X = U \oplus U' = V \oplus V'$. Then
$U \oplus V \oplus (U' \cap V') = (U \cap V') \oplus (V \cap U')
\oplus (U' \cap V') = X$, proving $U \oplus V \in \GrassB{X}$.
The second claim follows from $X = (U \cap V) \oplus (U \cap V')
\oplus (V \cap U') \oplus (U' \cap V')$.
\end{proof}

\begin{theorem}[Immersion theorem, {\cite[Theorem~4.15]{FHN2019}}]
\label{thm:tucker_immersion}
Assume \condINJ\ holds. Then for each $\bv \in \Tucker{\fr}{\VD}$:
\begin{equation}
\label{eq:ZD_split}
    \ZD{\bv} \in \GrassB{\VDnorm}.
\end{equation}
Consequently, $\Tucker{\fr}{\VD}$ is an immersed submanifold of
$\VDnorm$.
\end{theorem}

\begin{proof}
For each $\alpha \in D$: $\Wmin{\alpha}{\bv} \in
\GrassB{V_{\alpha,\|\cdot\|_\alpha}}$ and $\Umin{D\setminus\{\alpha\}}{\bv}$
is finite-dimensional. By Lemma~\ref{lem:two_faces},
$\Wmin{\alpha}{\bv} \atens \Umin{D\setminus\{\alpha\}}{\bv}
\in \GrassB{\VDnorm}$.
Also, ${}_a\bigotimes_\alpha \Umin{\alpha}{\bv} \in \GrassB{\VDnorm}$
(finite-dimensional, hence split by
Proposition~\ref{prop:finitedim_split}).
The $d+1$ summands in~\eqref{eq:ZD} have pairwise trivial intersection
(by the linear independence properties of minimal subspaces
and the direct sum structure), so Lemma~\ref{lem:direct_sum_G}
applied iteratively gives $\ZD{\bv} \in \GrassB{\VDnorm}$.
\end{proof}

\begin{corollary}[Best projection onto tangent space]
\label{cor:best_proj_tangent}
If $\VDnorm$ is additionally reflexive, then for any $\bv \in
\Tucker{\fr}{\VD}$ and $\dot{\bu} \in \VDnorm$, there exists
$\dot{\bv}_\mathrm{best} \in \ZD{\bv}$ such that
\begin{equation*}
    \|\dot{\bu} - \dot{\bv}_\mathrm{best}\|_D
    = \min_{\dot{\bv} \in \ZD{\bv}} \|\dot{\bu} - \dot{\bv}\|_D.
\end{equation*}
\end{corollary}

\begin{proof}
Since $\VDnorm$ is reflexive, every closed subspace is proximinal
(proximinal = best approximations exist). $\ZD{\bv}$ is closed
(it is a split subspace by Theorem~\ref{thm:tucker_immersion}).
\end{proof}

%------------------------------------------------------------
\section{The tensor space as a Banach manifold}
\label{sec:tucker_manifold:full}
%------------------------------------------------------------

\begin{corollary}[{\cite[Corollary~3.19]{FHN2019}}]
\label{cor:VD_manifold}
Under the continuity assumption on the tensor product map, the algebraic
tensor space $\VD$ itself is a $\mathcal{C}^\infty$-Banach manifold
(not modelled on a single Banach space), with atlas given by the
disjoint union of the Tucker manifold atlases over all $\fr \in \AD{\VD}$:
\begin{equation*}
    \VD = \bigsqcup_{\fr \in \AD{\VD}} \Tucker{\fr}{\VD}.
\end{equation*}
\end{corollary}

%------------------------------------------------------------
\section{Historical notes and comparison}
\label{sec:tucker_manifold:history}
%------------------------------------------------------------

The $\mathcal{C}^\infty$-Banach manifold structure of
$\Tucker{\fr}{\VD}$ in infinite-dimensional spaces was established
in~\cite{FHN2019}. Prior work had treated only finite-dimensional or
Hilbert space settings:

Koch and Lubich~\cite{KochLubich2007} proved that the Tucker manifold
is a quotient manifold in finite dimensions and used it for the
Dirac--Frenkel variational principle for the MCTDH method. Their
approach, however, uses the Hilbert structure (metric projection)
and does not extend to Banach spaces.

Uschmajew and Vandereycken~\cite{Uschmajew2012} studied the HT
manifold in finite dimensions as a quotient manifold
$\GL{r_1} \times \cdots \times \GL{r_d} \backslash \FTv / G$
for a gauge group $G$. The description in~\cite{FHN2019} is
fundamentally different: it provides an \emph{explicit atlas}
with local charts given by formula~\eqref{eq:local_rep}, valid
in the infinite-dimensional Banach setting.

The key novelties of~\cite{FHN2019} are:
\begin{enumerate}
\item The chart maps $\chartmap{\bv}$ are given explicitly in terms
    of the Grassmann chart maps $\Psi^{(\alpha)}_\bv$, making the
    geometric structure completely transparent.
\item The immersion theorem (Theorem~\ref{thm:tucker_immersion})
    holds in general Banach spaces under condition~\condINJ, without
    requiring a Hilbert or reflexive structure.
\item The fibre bundle structure identifies the typical fibre
    as $\RR^{\times r_\alpha}_*$, a finite-dimensional
    smooth manifold (an open subset of a Euclidean space).
\end{enumerate}

